%Abstract Page

\begin{titlepage}
\begin{center}

\vspace*{5\baselineskip}
\textbf{\large Abstract}

Mapping Initial-State Energy Distributions in Au+Au and p+p Collisions with sPHENIX

Emma McLaughlin
\end{center}
\begin{flushleft}
\hspace{10mm}This thesis presents the development of a calorimeter-based physics program for the sPHENIX detector at the Relativistic Heavy Ion Collider (RHIC). It first presents a viewpoint into the process of making calorimeter physics measurements including detector commissioning, reconstruction, and calibrations, and then details two physics measurements using the calorimeter detectors during RHIC Run 2024.
The transverse energy per unit pseudorapidity (\detdeta) is measured in Au+Au collisions at \snn = 200 GeV using the full sPHENIX calorimeter system over $|\eta| < 1.1$. This measurement is performed independently with the electromagnetic and hadronic calorimeters and good agreement between the two detector systems is found despite their differing calibration procedures. Results are reported differentially in centrality over the range $0$--$70$\% and as a function of the mean number of nucleon participant pairs (\Npart/2). Results are compared to previous measurements at RHIC and show improved $\eta$ granularity across the full centrality range and reduced uncertainties in peripheral events. This measurement of \detdeta establishes the performance of the sPHENIX calorimeter over more than an order of magnitude in dynamic range, demonstrating a strong understanding of the sPHENIX calorimeter response and establishing a foundation for the detector's jet physics program at RHIC.
The underlying event (UE) is measured as a function of leading jet $p_{T}$ in p+p collisions at \sqs = 200 GeV, utilizing calorimeter topoclusters to characterize soft particle production in regions azimuthally transverse to the hard scattering. The UE activity, given by the transverse region average transverse energy density, $\langle \Sigma E_T / \delta\eta\delta\phi \rangle$, is presented for both inclusive jet and dijet event selections. A consistent decrease in transverse energy density with increasing jet $p_{T}$ is observed in both selections, showing a clear anti-correlation between the hard scattering scale and soft particle production at RHIC energies. Comparisons to \pythia and \herwig are also included to highlight the sensitivity of the UE to the modeling of the interplay between soft particle production and the hard scattering. The origin of hard-soft correlation observed in the UE measurement deserves further investigation and possible mechanisms from studies in small systems like proton-ion collisions are presented. 

\end{flushleft}
\vspace*{\fill}
\end{titlepage}
